\chapter{Conclusions and Future Work}

\section{What the Campaign Established}

The campaign began with a broad wavefront-correction ambition and ended with a
more precise detector-side conclusion.

\begin{enumerate}
  \item A static phase-only D2NN preserves field-level invariants in the sampled
  unitary model. It is therefore not a random-turbulence wavefront corrector in
  the usual sense.
  \item Intensity observables such as PIB are not unitary invariants. They can
  improve, but those improvements must be interpreted on the correct plane.
  \item Normalized PIB alone is a misleading detector objective. The decisive
  detector-side quantity is absolute received power.
  \item Throughput must be built directly into the loss. Once that is done,
  static spatial D2NNs can produce modest but real received-power gains.
  \item Those gains are regime-dependent: weak and strong turbulence benefit
  differently, and distance-specific models transfer poorly across regimes.
  \item The corrected 4f FD2NN sweep showed only a near-tie, not a meaningful
  detector-side breakthrough.
\end{enumerate}

\section{Interpreting the Architecture Split}

The report does not support the claim that static diffractive optics are useless
for turbulent receivers. It supports a narrower architectural conclusion:
detector-centric spatial shaping is substantially more promising than static
Fourier-plane phase correction in the geometries explored here.

\section{MPLC as Context, Not a Verdict}

The committed MPLC comparison assets are useful as context because they remind
us that unitary mode-conversion optics can be engineered deliberately. They are
not a controlled head-to-head result for this campaign. We therefore use MPLC as
a future-work comparison target rather than as a verdict against D2NN.

\section{Future Work}

The most relevant follow-up directions are:
\begin{itemize}
  \item adaptive optics or adaptive D2NN front ends that can respond to
  realization-level turbulence changes;
  \item a controlled apples-to-apples comparison between spatial D2NN and 4f
  FD2NN under identical data, detector metrics, and throughput accounting;
  \item broader detector objectives, including multi-bucket and fiber-coupling
  surrogates;
  \item a cleaner comparison between passive D2NN front ends and MPLC-style
  mode-conversion systems.
\end{itemize}

\begin{breakthrough}
The central breakthrough of the campaign was conceptual rather than architectural:
once the objective was rewritten in terms of detector-side received power, the
experimental story became coherent. The successful branch was not ``make the
beam look sharp.'' It was ``deliver more useful energy to the detector.''
\end{breakthrough}
